\documentclass[10pt,a4paper]{amsart}
\usepackage{amsmath,amssymb,amsthm,graphicx}
\usepackage{verbatim}
\newenvironment{reference}{\comment}{\endcomment}
%\newenvironment{reference}{}{}
\newenvironment{slogan}{\comment}{\endcomment}
\newenvironment{history}{\comment}{\endcomment}

% For commutative diagrams we use Xy-pic
\usepackage[all]{xy}

% We use 2cell for 2-commutative diagrams.
\xyoption{2cell}
% Theorem environments.
%
\theoremstyle{plain}
\newtheorem{theorem}[subsection]{Theorem}
\newtheorem{proposition}[subsection]{Proposition}
\newtheorem{lemma}[subsection]{Lemma}

\theoremstyle{definition}
\newtheorem{definition}[subsection]{Definition}
\newtheorem{example}[subsection]{Example}
\newtheorem{exercise}[subsection]{Exercise}
\newtheorem{situation}[subsection]{Situation}

\theoremstyle{remark}
\newtheorem{remark}[subsection]{Remark}
\newtheorem{remarks}[subsection]{Remarks}

\numberwithin{equation}{subsection}

% Macros
%
\def\lim{\mathop{\mathrm{lim}}\nolimits}
\def\colim{\mathop{\mathrm{colim}}\nolimits}
\def\Spec{\mathop{\mathrm{Spec}}}
\def\Hom{\mathop{\mathrm{Hom}}\nolimits}
\def\Ext{\mathop{\mathrm{Ext}}\nolimits}
\def\SheafHom{\mathop{\mathcal{H}\!\mathit{om}}\nolimits}
\def\SheafExt{\mathop{\mathcal{E}\!\mathit{xt}}\nolimits}
\def\Sch{\mathit{Sch}}
\def\Mor{\mathop{\mathrm{Mor}}\nolimits}
\def\Ob{\mathop{\mathrm{Ob}}\nolimits}
\def\Sh{\mathop{\mathit{Sh}}\nolimits}
\def\NL{\mathop{N\!L}\nolimits}
\def\CH{\mathop{\mathrm{CH}}\nolimits}
\def\proetale{{pro\text{-}\acute{e}tale}}
\def\etale{{\acute{e}tale}}
\def\QCoh{\mathit{QCoh}}
\def\Ker{\mathop{\mathrm{Ker}}}
\def\Im{\mathop{\mathrm{Im}}}
\def\Coker{\mathop{\mathrm{Coker}}}
\def\Coim{\mathop{\mathrm{Coim}}}

% Boxtimes
%
\DeclareMathSymbol{\boxtimes}{\mathbin}{AMSa}{"02}

%
% Macros for moduli stacks/spaces
%
\def\QCohstack{\mathcal{QC}\!\mathit{oh}}
\def\Cohstack{\mathcal{C}\!\mathit{oh}}
\def\Spacesstack{\mathcal{S}\!\mathit{paces}}
\def\Quotfunctor{\mathrm{Quot}}
\def\Hilbfunctor{\mathrm{Hilb}}
\def\Curvesstack{\mathcal{C}\!\mathit{urves}}
\def\Polarizedstack{\mathcal{P}\!\mathit{olarized}}
\def\Complexesstack{\mathcal{C}\!\mathit{omplexes}}
% \Pic is the operator that assigns to X its picard group, usage \Pic(X)
% \Picardstack_{X/B} denotes the Picard stack of X over B
% \Picardfunctor_{X/B} denotes the Picard functor of X over B
\def\Pic{\mathop{\mathrm{Pic}}\nolimits}
\def\Picardstack{\mathcal{P}\!\mathit{ic}}
\def\Picardfunctor{\mathrm{Pic}}
\def\Deformationcategory{\mathcal{D}\!\mathit{ef}}
\usepackage{cleveref}
\crefname{originallemma}{original lemma}{original lemmas}
\Crefname{originallemma}{Original lemma}{Original lemmas}
\makeatletter
\newlabel{chow-lemma-intersect-properly}{{0FFD}{source}{}{sourcepointer-1}{}}
\newlabel{chow-lemma-intersect-properly@cref}{{[originallemma][0][]0FFD}{[1][0][]source}{}{}{}}
\makeatother
\crefname{originallemma}{original lemma}{original lemmas}
\Crefname{originallemma}{Original lemma}{Original lemmas}
\makeatletter
\newlabel{morphisms-lemma-finite-type-ample-very-ample}{{01VU}{source}{}{sourcepointer-2}{}}
\newlabel{morphisms-lemma-finite-type-ample-very-ample@cref}{{[originallemma][0][]01VU}{[1][0][]source}{}{}{}}
\makeatother
\crefname{originalproposition}{original proposition}{original propositions}
\Crefname{originalproposition}{Original proposition}{Original propositions}
\makeatletter
\newlabel{properties-proposition-characterize-ample}{{01Q3}{source}{}{sourcepointer-3}{}}
\newlabel{properties-proposition-characterize-ample@cref}{{[originalproposition][0][]01Q3}{[1][0][]source}{}{}{}}
\makeatother
\crefname{originallemma}{original lemma}{original lemmas}
\Crefname{originallemma}{Original lemma}{Original lemmas}
\makeatletter
\newlabel{varieties-lemma-baby-stein}{{0FD1}{source}{}{sourcepointer-4}{}}
\newlabel{varieties-lemma-baby-stein@cref}{{[originallemma][0][]0FD1}{[1][0][]source}{}{}{}}
\makeatother
\crefname{originallemma}{original lemma}{original lemmas}
\Crefname{originallemma}{Original lemma}{Original lemmas}
\makeatletter
\newlabel{varieties-lemma-bertini}{{0FD6}{source}{}{sourcepointer-5}{}}
\newlabel{varieties-lemma-bertini@cref}{{[originallemma][0][]0FD6}{[1][0][]source}{}{}{}}
\makeatother
\crefname{originallemma}{original lemma}{original lemmas}
\Crefname{originallemma}{Original lemma}{Original lemmas}
\makeatletter
\newlabel{varieties-lemma-generate-over-complement}{{0FD5}{source}{}{sourcepointer-6}{}}
\newlabel{varieties-lemma-generate-over-complement@cref}{{[originallemma][0][]0FD5}{[1][0][]source}{}{}{}}
\makeatother
\crefname{originallemma}{original lemma}{original lemmas}
\Crefname{originallemma}{Original lemma}{Original lemmas}
\makeatletter
\newlabel{varieties-lemma-proper-geometrically-reduced-global-sections}{{0BUG}{source}{}{sourcepointer-7}{}}
\newlabel{varieties-lemma-proper-geometrically-reduced-global-sections@cref}{{[originallemma][0][]0BUG}{[1][0][]source}{}{}{}}
\makeatother
\crefname{originallemma}{original lemma}{original lemmas}
\Crefname{originallemma}{Original lemma}{Original lemmas}
\makeatletter
\newlabel{varieties-lemma-smooth-geometrically-normal}{{056T}{source}{}{sourcepointer-8}{}}
\newlabel{varieties-lemma-smooth-geometrically-normal@cref}{{[originallemma][0][]056T}{[1][0][]source}{}{}{}}
\makeatother
\title{Zero cycles are generated by classes of points with separable residue field}
\author{The Stacks Project Authors}
\date{stacks-project:weil.tex:local-snapshot-20261012; GNU FDL 1.2 or later}
\begin{document}
\maketitle
\noindent\textit{Editorial note.} All delivered mathematical text is the author's own, copied from the pinned Stacks Project chapter file weil.tex; no statement or proof step was rewritten, reordered or removed.\par\medskip
\noindent\textit{Editorial note.} The counted proof is the author's complete argument: Prefilter hit re-read per the 2026-10-14 12:4x THIN-criterion patch (pb>=800 / pw>=100): the proof body was read in full at source and its main line is complete; the flagged phrase is incidental and is quoted verbatim below. Verbatim, proof line 2295: '(details omitted)' after the citation of the Chow Homology lemma on intersecting cycles that meet properly; that application is a routine verification and is not load-bearing for the construction of the complete intersection.\par\medskip
\noindent\textit{Editorial note.} The label lemma-generated-by-separable names the counted statement, and the proof environment immediately following it in the source is delivered verbatim as its proof.\par\medskip
\noindent\textit{Editorial note.} The counted statement is carried by the author's own lemma environment in the chapter source.\par\medskip
\noindent\textit{Editorial note.} Prefilter hit re-read per the 2026-10-14 12:4x THIN-criterion patch: proof body (proof\_lines 2239-2310, pb=3023, pw=461) read in full; main line complete.\par\medskip
\noindent\textit{Editorial note.} Editorial note (verbatim): 'in \$\textbackslash{}CH\_0(X)\$ by repeated application of Chow Homology, Lemma ... (details omitted).'  -  the omitted detail is the routine application of the cited intersection lemma.\par\medskip
\noindent\textit{Editorial note.} Pre-open gate: python operations/agent-authoring/tools/precard\_gate.py --source weil.tex --label lemma-generated-by-separable --start 2232 --end 2310 --local-id 28163 returned RESULT=OK. Pre-assembly tag self-check: canonical tag resolved and not present in reference/public-items.jsonl.\par\medskip
\noindent\textit{Editorial note.} The statement and proof are copied from the pinned chapter file; the cited inputs are kept as external pointers with their canonical Stacks tags.\par\medskip
\noindent\textit{Allowed standard prerequisites (editorial).}\begin{itemize}
\item Structure of the ring of global sections of a proper geometrically reduced scheme. Varieties, the lemmas on smooth geometrically normal schemes, on global sections of proper geometrically reduced schemes, and the Stein factorization statement for varieties. Reduces to a geometrically integral scheme over a finite separable extension of the base field.
\item Ampleness, generation over the complement of a point and Bertini type statements. Morphisms, the existence of very ample twists, Varieties, the lemma on generating a sheaf over the complement of a point, Properties, the characterization of ample invertible sheaves, and Varieties, the Bertini statement. Produces sections whose vanishing loci meet properly in the chosen point and are otherwise smooth, giving an etale zero dimensional residual scheme.
\item Intersecting divisors that meet properly. Chow Homology, the lemma on the intersection product of cycles meeting properly. Identifies the product of the classes of the vanishing loci with the sum of the class of the point and the class of the residual scheme.
\end{itemize}\medskip
\ifdefined\hypertarget\hypertarget{sourcepointer-1}{}\fi
\noindent\textit{Source reference.} Original lemma 0FFD: sources/stacks-chow/tags.txt, source lines 19283--19283. Prerequisite pointer invoked at proof line(s) 2304; it names a standard fact of a sibling chapter that the counted argument uses, and every step of the counted proof built on it is retained verbatim in the delivered body.\par\medskip
\ifdefined\hypertarget\hypertarget{sourcepointer-2}{}\fi
\noindent\textit{Source reference.} Original lemma 01VU: sources/stacks-chow/tags.txt, source lines 2459--2459. Prerequisite pointer invoked at proof line(s) 2267; it names a standard fact of a sibling chapter that the counted argument uses, and every step of the counted proof built on it is retained verbatim in the delivered body.\par\medskip
\ifdefined\hypertarget\hypertarget{sourcepointer-3}{}\fi
\noindent\textit{Source reference.} Original proposition 01Q3: sources/stacks-chow/tags.txt, source lines 2252--2252. Prerequisite pointer invoked at proof line(s) 2269; it names a standard fact of a sibling chapter that the counted argument uses, and every step of the counted proof built on it is retained verbatim in the delivered body.\par\medskip
\ifdefined\hypertarget\hypertarget{sourcepointer-4}{}\fi
\noindent\textit{Source reference.} Original lemma 0FD1: sources/stacks-chow/tags.txt, source lines 19201--19201. Prerequisite pointer invoked at proof line(s) 2249; it names a standard fact of a sibling chapter that the counted argument uses, and every step of the counted proof built on it is retained verbatim in the delivered body.\par\medskip
\ifdefined\hypertarget\hypertarget{sourcepointer-5}{}\fi
\noindent\textit{Source reference.} Original lemma 0FD6: sources/stacks-chow/tags.txt, source lines 19206--19206. Prerequisite pointer invoked at proof line(s) 2290; it names a standard fact of a sibling chapter that the counted argument uses, and every step of the counted proof built on it is retained verbatim in the delivered body.\par\medskip
\ifdefined\hypertarget\hypertarget{sourcepointer-6}{}\fi
\noindent\textit{Source reference.} Original lemma 0FD5: sources/stacks-chow/tags.txt, source lines 19205--19205. Prerequisite pointer invoked at proof line(s) 2268; it names a standard fact of a sibling chapter that the counted argument uses, and every step of the counted proof built on it is retained verbatim in the delivered body.\par\medskip
\ifdefined\hypertarget\hypertarget{sourcepointer-7}{}\fi
\noindent\textit{Source reference.} Original lemma 0BUG: sources/stacks-chow/tags.txt, source lines 14874--14874. Prerequisite pointer invoked at proof line(s) 2248; it names a standard fact of a sibling chapter that the counted argument uses, and every step of the counted proof built on it is retained verbatim in the delivered body.\par\medskip
\ifdefined\hypertarget\hypertarget{sourcepointer-8}{}\fi
\noindent\textit{Source reference.} Original lemma 056T: sources/stacks-chow/tags.txt, source lines 6721--6721. Prerequisite pointer invoked at proof line(s) 2247; it names a standard fact of a sibling chapter that the counted argument uses, and every step of the counted proof built on it is retained verbatim in the delivered body.\par\medskip
\begin{lemma}
\label{lemma-generated-by-separable}
Let $k$ be a field. Let $X$ be a smooth projective scheme over $k$.
Then $\CH_0(X)$ is generated by classes of closed points whose residue
fields are separable over $k$.
\end{lemma}

\begin{proof}
The lemma is immediate if $k$ has characteristic $0$ or is perfect.
Thus we may assume $k$ is an infinite field of characteristic $p > 0$.

\medskip\noindent
We may assume $X$ is irreducible of dimension $d$.
Then $k' = H^0(X, \mathcal{O}_X)$ is a finite separable field
extension of $k$ and that $X$ is geometrically integral over $k'$.
See Varieties, Lemmas \ref{varieties-lemma-smooth-geometrically-normal},
\ref{varieties-lemma-proper-geometrically-reduced-global-sections}, and
\ref{varieties-lemma-baby-stein}. We may and do replace $k$ by $k'$
and assume that $X$ is geometrically integral.

\medskip\noindent
Let $x \in X$ be a closed point. To prove the lemma we are going to show that
$[x] \in \CH_0(X)$ is rationally equivalent to an integer linear
combination of classes of closed points whose residue fields
are separable over $k$. Choose an ample invertible
$\mathcal{O}_X$-module $\mathcal{L}$. Set
$$
V = \{s \in H^0(X, \mathcal{L}) \mid s(x) = 0 \}
$$
After replacing $\mathcal{L}$ by a power we may assume
(a) $\mathcal{L}$ is very ample, (b) $V$ generates
$\mathcal{L}$ over $X \setminus x$, (c) the morphism
$X \setminus x \to \mathbf{P}(V)$ is an immersion, (d)
the map $V \to \mathfrak m_x\mathcal{L}_x/\mathfrak m_x^2\mathcal{L}_x$
is surjective, see
Morphisms, Lemma \ref{morphisms-lemma-finite-type-ample-very-ample},
Varieties, Lemma \ref{varieties-lemma-generate-over-complement}, and
Properties, Proposition \ref{properties-proposition-characterize-ample}.
Consider the set
$$
V^d \supset U =
\{
(s_1, \ldots, s_d) \in V^d \mid s_1, \ldots, s_d
\text{ generate }
\mathfrak m_x\mathcal{L}_x/\mathfrak m_x^2\mathcal{L}_x
\text{ over }\kappa(x)
\}
$$
Since $\mathcal{O}_{X, x}$ is a regular local ring of dimension $d$
we have $\dim_{\kappa(x)}(\mathfrak m_x/\mathfrak m_x^2) = d$
and hence we see that $U$ is a nonempty (Zariski) open of $V^d$.
For $(s_1, \ldots, s_d) \in U$ set $H_i = Z(s_i)$. Since
$s_1, \ldots, s_d$ generate $\mathfrak m_x\mathcal{L}_x$
we see that
$$
H_1 \cap \ldots \cap H_d = x \amalg Z
$$
scheme theoretically for some closed subscheme $Z \subset X$.
By Bertini (in the form of Varieties, Lemma \ref{varieties-lemma-bertini})
for a general element $s_1 \in V$ the scheme $H_1 \cap (X \setminus x)$
is smooth over $k$ of dimension $d - 1$.
Having chosen $s_1$, for a general element
$s_2 \in V$ the scheme $H_1 \cap H_2 \cap (X \setminus x)$
is smooth over $k$ of dimension $d - 2$. And so on.
We conclude that for sufficiently general
$(s_1, \ldots, s_d) \in U$ the scheme $Z$ is \'etale over $\Spec(k)$.
In particular $H_1 \cap \ldots \cap H_d$ has dimension $0$
and hence
$$
[H_1] \cdot \ldots \cdot [H_d] = [x] + [Z]
$$
in $\CH_0(X)$ by repeated application of
Chow Homology, Lemma \ref{chow-lemma-intersect-properly} (details omitted).
This finishes the proof as it shows that $[x] \sim_{rat} - [Z] + [Z']$
where $Z' = H'_1 \cap \ldots \cap H'_d$ is a general complete
intersection of vanishing loci of sufficiently general sections
of $\mathcal{L}$ which will be \'etale over $k$ by the same argument
as before.
\end{proof}

\end{document}
